\chapter{Method}
\label{ch:method}

\section{What one run is}

An evaluation run takes one split, one engine and one seed, loads every form in
that split through a real browser, extracts descriptors, classifies every keyed
field, runs the audit engine over the same page, and writes five files: a run
log with one row per classified field, a manifest, a metrics document, a
confusion document, and a findings file.

The run log row carries the form identifier, the template, the locale, the tier,
the field selector, the answer key label, the predicted label, the confidence,
the kind of confidence, the runner up, the declared token, the per field
classification latency, and the engine's own name. The manifest carries the
schema version, the commit and whether the tree was dirty, resolved dependency
versions, the corpus manifest digest, the split file digest, the seed, the
engine's self description, the command line, a machine descriptor and the start
and end timestamps.

Three of those deserve a sentence each.

\textbf{The engine field carries the engine's own name.} The specification uses
one label in its run log schema and a different one in its headline table for
the same engine. Carrying two names for one engine across a repository is how a
comparison table becomes ambiguous, so the run log records what the engine calls
itself and a report labels its column however it likes.

\textbf{The confidence kind is mapped rather than copied.} The schema's
vocabulary and the engines' own words are not the same strings, because the
engines' words were made load bearing in the renderers first. The mapping is
written down once and is total: an engine whose word is not in it raises rather
than defaulting to something plausible.

\textbf{The findings file is a fifth artefact the specification does not list.}
Whether a page \emph{needed} an accusation depends on the difference between
declaring nothing, declaring off, and declaring a token outside the
specification, and the run log row carries a declared token that is null in all
three cases. Without the extra file the finding level significance test could not
be reproduced from committed artefacts, and a number nobody can recompute is a
number law 3 does not allow into a report.

\section{Metrics}

Macro-F1 and micro-F1 over the union of observed labels, per slice; per label
precision, recall and F1; a confusion table; abstention rate and accuracy on the
fields the engine did commit to; finding level precision and recall per finding
code; and per field latency percentiles beside the wall clock breakdown.

\textbf{Both averages are always reported.} They disagree here about which of
the two classical engines leads, and a table carrying one of them would be
picking a winner by picking a denominator. The mechanism of the disagreement is
the abstention asymmetry and it is explained where the numbers are, in
Chapter~\ref{ch:headline}.

\textbf{A cell is reported only when it holds enough fields.} Below the
reporting minimum a cell renders the words \emph{insufficient data}, not a
blank, because a blank reads as an oversight and this is a decision. The
threshold is a single constant in the metrics module, imported by everything
that reports, with no per table override, and it was chosen before any corpus
had been generated and before any accuracy had been computed, which is the only
point at which it could be chosen without knowledge of its effect. On this
benchmark it fires on exactly one thing, and what it fires on is itself a
finding: see Section~\ref{sec:accusations}.

\section{The statistical procedure}

The unit of resampling is the template, because fields inside one template share
an author. The test is a paired sign flip permutation at the cluster level,
exhaustive when the arrangement space is small enough to enumerate, with
Benjamini and Hochberg step up correction across the whole comparison family
\cite{benjamini1995}, and a practical effect gate expressed as an absolute
difference in the metric.

\begin{table}[htbp]
  \centering
  \caption{The pre-registered statistical policy and the power the realised
  design has. Every value in the upper half was fixed before the runs.}
  \label{tab:design}
  \begingroup\small\setlength{\tabcolsep}{4pt}
\begin{tabular}{lr}
\toprule
\textbf{Design quantity} & \textbf{Value} \\
\midrule
Resampling unit & \texttt{template\textbackslash\{\}\_id} \\
Clusters in the test split & \texttt{10} \\  % results: experiments/results/analysis/2026-08-26T21-20-11Z_p6-analysis_6457ac7/analysis.json
Sign flip arrangements & \texttt{1024} \\  % results: experiments/results/analysis/2026-08-26T21-20-11Z_p6-analysis_6457ac7/analysis.json
Smallest attainable two sided p & \texttt{0.001953} \\  % results: experiments/results/analysis/2026-08-26T21-20-11Z_p6-analysis_6457ac7/analysis.json
Exhaustive rather than sampled & \texttt{yes} \\
Alpha & \texttt{0.050000} \\  % results: experiments/results/analysis/2026-08-26T21-20-11Z_p6-analysis_6457ac7/analysis.json
False discovery rate & \texttt{0.050000} \\  % results: experiments/results/analysis/2026-08-26T21-20-11Z_p6-analysis_6457ac7/analysis.json
Practical threshold, absolute & \texttt{0.020000} \\  % results: experiments/results/analysis/2026-08-26T21-20-11Z_p6-analysis_6457ac7/analysis.json
Comparison family size & \texttt{33} \\  % results: experiments/results/analysis/2026-08-26T21-20-11Z_p6-analysis_6457ac7/analysis.json
Seed & \texttt{20260825} \\  % results: experiments/results/analysis/2026-08-26T21-20-11Z_p6-analysis_6457ac7/analysis.json
\bottomrule
\end{tabular}
\endgroup
\par\smallskip\noindent{\footnotesize Source: \rpath{experiments/results/analysis/2026-08-26T21-20-11Z_p6-analysis_6457ac7/analysis.json}}

\end{table}

Two decisions in Table~\ref{tab:design} were fixed while fixing them was still
costless, and both would have been contaminated by seeing the numbers.

\textbf{The comparison family is all three engine pairs over the same eleven
comparisons, corrected together.} A three engine benchmark could reasonably have
corrected across one pair, or two, or all three. The temptation is specific and
foreseeable: the comparison this project exists to make is the one against the
language model, and a family containing only that pair would be a third the size
and would produce smaller adjusted p values for exactly the comparisons the
project most wants to be able to report. The size was therefore fixed in the
prediction file, along with the stated consequence that every comparison between
the two classical engines would carry a larger adjusted p than the same
comparison had carried when it was corrected inside a smaller family. It did, and
nobody can report that as a change in the underlying result, because the file
said it would happen.

\textbf{The practical effect gate is absolute.} The external harness's gate is a
relative percentage; this project registered an absolute difference in the
metric. A relative gate turns one pre-registered rule into a different rule in
every slice, because it means something different against a large baseline than
against a small one.

\section{The dependency boundary, and what it cost}

Permutation testing, multiple comparison correction and the verdict vocabulary
are not implemented here. They live in a separate, independently released
evaluation harness \cite{triage} that this project consumes as an ordinary
dependency: not vendored, not submoduled, not copied file by file.

That separation is a deliberate engineering claim. A piece of infrastructure
with exactly one consumer has not been shown to be infrastructure; it has been
shown to be part of that one program. This project is the second consumer, across
a real package boundary, and the friction the boundary exposes is the actual
finding. Vendoring the code would have destroyed that evidence by making the
boundary unobservable.

The friction split cleanly, and the split is more interesting than either
outcome would have been alone.

\textbf{The statistical primitives fit exactly and were used unchanged.} The p
value estimator takes the null distribution as an argument, which is precisely
the seam a caller with its own resampling scheme needs. The correction procedure
needed nothing at all. Its verdict vocabulary already contained the three
categories required plus a fourth for a design whose smallest attainable p value
sits above alpha, which turned out to be the category every comparison in the
first analysis landed in.

\textbf{Everything shaped like a training run did not fit.} The ingestion layer
claims a run log of this shape and then refuses it for having no step field, and
there is no step to add, because the file has no time axis. No public entry point
accepts a cluster assignment. Neither comparison mode is paired, and both reduce
a metric series to a final window mean, which a categorical per field outcome
does not have.

Five issues are filed on that repository, each with a concrete API proposal, and
the ledger in this repository records each one with its status and the
workaround, if any, that stands here in the meantime. Two workarounds exist and
both are labelled in every result file they produce, because a workaround that is
invisible in the output is indistinguishable from a capability: the clustered
null is built here and handed to the harness's estimator, and the practical
effect gate is applied here because the harness's is relative.

The anticipated friction was aimed one layer too low. The ledger predicted,
before the integration began, that categorical outcome support in the permutation
machinery would be the problem. It was not needed at all: the estimator never
sees an outcome, only an effect and a null, both numbers. The categorical part of
the problem lives entirely in the statistic and the statistic is the caller's.
What actually broke was one layer above, in the data model.

One consequence of the boundary is stated here rather than discovered later. The
harness is a development dependency, not a runtime one, because installing it
pulls more than a dozen transitive packages including a training metrics logging
stack, and the audit path never imports it. \textbf{The shipped tool therefore
cannot compute its own significance tests.} The narrower extra that would fix it
is proposed in the ledger.

\section{Reproducibility, and the one artefact that is marked dirty}

One seed reaches the generator, the split assignment, the training script and
any sampling in evaluation, through explicitly passed generator objects. The lock
file is committed and continuous integration installs from it. Every result
directory carries a manifest. Model artefacts are content addressed into every
run that used them. Nothing is ever regenerated to fix a number: if a result is
wrong, a new run is made with a new identifier and the old one stays, because
result files are append only history.

A run produced from a working tree with uncommitted changes is marked dirty in
its manifest and is not citable, and the traceability job fails on a citation of
one. Every evaluation result cited anywhere in this report is marked clean.

\textbf{One artefact is not, and it is the training manifest of the shipped
model.} The model was fitted in the same working session that authored the new
corpus templates, so the tree carried uncommitted changes at the moment of
fitting, and the manifest says so honestly rather than hiding it. What follows
from that is worth being precise about rather than either ignoring or
overstating:

\begin{itemize}
  \item The model artefacts themselves are committed and are content addressed
    into every evaluation run manifest by digest, so \emph{which} model produced
    the headline numbers is not in doubt.
  \item The corpus is content addressed the same way, and its digest in the
    training manifest matches the digest in every evaluation manifest, so the
    model was fitted on the same corpus it was evaluated against.
  \item What is genuinely weaker is the development split evidence: the numbers
    in Tables~\ref{tab:training}, \ref{tab:sweep}, \ref{tab:calibration} and
    \ref{tab:parity} come from a run whose tree state was not pinned to a
    commit. They are reported because they describe how the shipped model was
    built and they should not be silently dropped, and they are flagged here
    because a reader is entitled to know that they sit one notch below the test
    split numbers on the same scale the project applies to everything else.
  \item The fix is procedural rather than technical, it is not applied
    retroactively because that would be a regeneration, and it is listed in
    Chapter~\ref{ch:future}: commit the corpus change first, then train from the
    clean tree.
\end{itemize}
